% ============================================================================
%  SECCIÓN 3.17: COMPARACIÓN DEL PDF 1 CON EL PDF 2 (PASO 19)
% ============================================================================

\section{Comparaci\'on del PDF 1 con el PDF 2}

Esta secci\'on es fundamental para evidenciar y contrastar el avance
metodol\'ogico y t\'ecnico del proyecto entre la primera entrega (PDF 1) y la
segunda entrega (PDF 2). A continuaci\'on, se presenta la tabla comparativa que
sintetiza la evoluci\'on de los elementos clave de la propuesta y su estado
actual de implementaci\'on.

\begin{table}[H]
  \centering
  \caption{Comparativa de avances entre el PDF 1 y el PDF 2}
  \label{tab:comparacion-pdf1-pdf2}
  \contenidoConFuente{%
    \begin{tablaAPA}
      \begin{tabular}{@{}p{3.2cm}p{5.8cm}p{6.6cm}@{}}
        \toprule
        \textbf{Elemento} & \textbf{PDF 1 (Entrega Inicial)} & \textbf{PDF 2 (Avance Actual)} \\
        \midrule
        Nombre del proyecto & A\'un no definido & \nombreApp{} (denominaci\'on provisional) \\
        \addlinespace
        Problema & Definido & Se mantiene (ratificado) \\
        \addlinespace
        Usuarios & Identificados & Se mantienen (clasificados por roles) \\
        \addlinespace
        Funcionalidades & Definidas & Definidas y por desarrollar (priorizadas en MVP) \\
        \addlinespace
        Wireframes & Dise\~nados & Dise\~nados y por implementar (primera pantalla implementada) \\
        \addlinespace
        Tecnolog\'ia & Propuesta & Propuesta, configurada y por ser testeadas \\
        \addlinespace
        GitHub & No establecido (opcional) & Repositorio establecido y configurado \\
        \bottomrule
      \end{tabular}
    \end{tablaAPA}%
  }{Elaboraci\'on propia en base a la matriz de comparaci\'on de entregas.}
\end{table}

Como se evidencia en la tabla comparativa, el proyecto ha transitado desde una
fase inicial de definici\'on conceptual hacia una etapa de estructuraci\'on
formal y preparaci\'on t\'ecnica, consolidando la identidad del sistema bajo el
nombre provisional \nombreApp{}, estableciendo el repositorio colaborativo en
GitHub y sentando las bases t\'ecnicas y de dise\~no para el desarrollo integral
de los m\'odulos funcionales.
